\documentclass[a4paper,12pt]{article}
\usepackage[utf8]{inputenc}
\usepackage[russian]{babel}
\usepackage{amsmath,amssymb}
\usepackage{graphicx}
\usepackage{geometry}
\usepackage{float}
\usepackage{hyperref}
\usepackage{array}
\usepackage{booktabs}
\usepackage{caption}

\geometry{left=2cm,right=2cm,top=2cm,bottom=2cm}

\begin{document}

\section*{Лабораторная работа №2}

\subsection*{Исходное регулярное выражение}
$(dcdc^{*}c(d^{*}|e^{*}e)|e(cc)^{*}(ee|d)|cac|dd)^{*}c(a|c|d|e)$

\subsection*{1. Минимальный ДКА}

\begin{figure}[H]
\centering
\includegraphics[width=0.9\textwidth]{1.png}
\caption{Диаграмма минимального ДКА}
\end{figure}

\textbf{Начальное состояние:} 0 \\
\textbf{Финальные состояния:} 15, 16, 17, 18, 19, 20, 21

\begin{table}[H]
\centering
\caption{Таблица переходов минимального ДКА}
\begin{tabular}{|c|c|c|c|c|c|}
\hline
Состояние & a & b & c & d & e \\
\hline
0 & 7 & 7 & 11 & 4 & 5 \\
1 & 7 & 7 & 7 & 6 & 7 \\
2 & 7 & 7 & 5 & 7 & 7 \\
3 & 7 & 7 & 7 & 7 & 0 \\
4 & 7 & 7 & 1 & 0 & 7 \\
5 & 7 & 7 & 2 & 0 & 3 \\
6 & 7 & 7 & 10 & 7 & 7 \\
7 & 7 & 7 & 7 & 7 & 7 \\
8 & 7 & 7 & 12 & 8 & 5 \\
9 & 7 & 7 & 13 & 8 & 9 \\
10 & 7 & 7 & 14 & 8 & 9 \\
11 & 17 & 7 & 15 & 15 & 15 \\
12 & 17 & 7 & 15 & 18 & 15 \\
13 & 17 & 7 & 16 & 15 & 15 \\
14 & 17 & 7 & 21 & 19 & 20 \\
15 & 7 & 7 & 7 & 7 & 7 \\
16 & 7 & 7 & 2 & 0 & 3 \\
17 & 7 & 7 & 0 & 7 & 7 \\
18 & 7 & 7 & 10 & 7 & 7 \\
19 & 7 & 7 & 12 & 8 & 5 \\
20 & 7 & 7 & 13 & 8 & 9 \\
21 & 17 & 7 & 21 & 19 & 20 \\
\hline
\end{tabular}
\end{table}

\textbf{Таблица классов эквивалентности:}

\href{https://github.com/TonitaN/FormalLanguageTheory/blob/main/2023/lect_tfl_3.pdf}{Согласно лекции строим}
\begin{table}[H]
\centering
\caption{Матрица эквивалентности (строки - классы, столбцы - суффиксы)}
\begin{tabular}{|c|*{10}{c|}}
\hline
Слова/Префиксы & $\varepsilon$ & a & cc & dcd & dc & eeedce & cdcd & cca & dcca & eca \\
\hline
$\varepsilon$ & - & - & + & - & - & - & - & - & - & - \\
dc & - & - & - & - & - & - & - & - & + & - \\
ec & - & - & - & - & - & - & + & - & - & - \\
ee & - & - & - & - & - & - & - & - & - & + \\
d & - & - & - & + & - & - & - & - & - & - \\
e & - & - & - & + & - & + & - & - & - & - \\
dcd & - & - & - & - & - & - & + & + & - & - \\
a & - & - & - & - & - & - & - & - & - & - \\
dcdcd & - & - & + & + & - & - & - & - & - & - \\
dcdce & - & - & + & + & - & + & - & - & - & + \\
dcdc & - & - & + & + & - & + & + & + & - & + \\
c & - & + & - & - & - & - & - & - & - & - \\
dcdcdc & - & + & - & - & - & - & - & - & + & - \\
dcdcec & - & + & - & - & - & - & + & - & - & - \\
dcdcc & - & + & + & + & - & + & + & + & - & + \\
cc & + & - & - & - & - & - & - & - & - & - \\
dcdcecc & + & - & - & + & - & + & - & - & - & - \\
ca & + & - & - & - & - & - & - & + & - & - \\
dcdcdcd & + & - & - & - & - & - & + & + & - & - \\
dcdccd & + & - & + & + & - & - & - & - & - & - \\
dcdcce & + & - & + & + & - & + & - & - & - & + \\
dcdccc & + & + & + & + & - & + & + & + & - & + \\
\hline
\end{tabular}
\end{table}
В данной таблице все строки описывают состояние, также они попарно различны, потому автомат минимален

\subsection*{2. Меньший НКА}
\begin{figure}[H]
\centering
\includegraphics[width=0.9\textwidth]{2.png}
\caption{Диаграмма НКА}
\end{figure}



\textbf{Таблица классов эквивалентности для минимального подмножества из трети состояний:}

\begin{table}[H]
\centering
\caption{Матрица эквивалентности (строки - слова, столбцы - префиксы)}
\begin{tabular}{|c|*{6}{c|}}
\hline
Слова/Префиксы & $\varepsilon$ & a & ca & dca & cdca & eca \\
\hline
ca & + & - & - & - & - & - \\
c & - & + & - & - & - & - \\
$\varepsilon$ & - & - & + & - & - & - \\
e & - & - & - & + & - & - \\
ec & - & - & - & - & + & - \\
ee & - & - & - & - & - & + \\
\hline
\end{tabular}
\end{table}
Здесь получаем верхне-треугольную матрицу, потому частично минимален автомат.

\subsection*{3. Переключающийся автомат}
Заметим, что у нас после $dcd$ у нас либо слово кончается, либо идет слово удовлетворяющее $cc^*(d^*|e^*e)$, а затем уже только может идти что угодно, тогда этот момент выкидываем из первоначального выражения, получая минимальный ДКА меньшего размера, далее пересекаем с автоматом, который проверяет вышеупомянутое свойство. Но состояний получится столько же. Можно было еще попробовать попересечь, изучая аналогично выражение, но автомат будет скорее лишь расти, особенно учитывая, что в алфавите целых 4 буквы, поэтому если я буду добавлять свойство, то итоговое число состояний будет лишь расти и автомат меньшего размера получить не получится. Того наш автомат будет по факту пересечением двух: 1. Где свойство даннное в выражении регулярном не учитывается, 2. Где это свойство строго проверяется, если свойство не выполняется, то уводим в выбранную ловушку.
\begin{figure}[H]
\centering
\includegraphics[width=0.8\textwidth]{3.png}
\caption{Диаграмма переключающегося автомата}
\label{fig:switching}
\end{figure}

\textbf{Таблица классов эквивалентности для переключающегося автомата:}

\begin{table}[H]
\centering
\caption{Матрица эквивалентности переключающегося автомата}
\begin{tabular}{|c|*{6}{c|}}
\hline
Слова/Префиксы & dc & a & eca & cc & dcd & $\varepsilon$ \\
\hline
dcdccc & - & + & + & + & + & + \\
dcdcce & - & - & + & + & + & + \\
dcdccd & - & - & - & + & + & + \\
dcdcecc & - & - & - & - & + & + \\
cc & - & - & - & - & - & + \\
$\varepsilon$ & - & - & - & + & - & - \\
\hline
\end{tabular}
\end{table}

Получаем верхне-треугольную матрицу с + вверху и - по диагонали.

\subsection*{4. Расширенное регулярное выражение}

Расширенное регулярное выражение имеет вид:

\begin{verbatim}
^(dcdc+d*|dcdc+e+|e(cc)*ee|e(cc)*d|cac|dd)*c.$
\end{verbatim}

Исходное регулярное выражение было преобразовано следующим образом:

\begin{enumerate}
\item \textbf{Последний символ}: В исходном выражении последняя буква могла быть любой из множества $(a|c|d|e)$, поэтому она заменена на метасимвол точки (.), который означает любой символ.

\item \textbf{Положительные итерации}: Введена операция положительной итерации (+), которая означает одно или более повторений:
\begin{itemize}
\item \texttt{dcdc+} заменяет \texttt{dcdc*c} - последовательность "dcdc" и одна или более букв 'c'
\item \texttt{e+} заменяет \texttt{e*e} - одна или более букв 'e'
\end{itemize}

\item \textbf{Ограничения на использование других операций}:
\begin{itemize}
\item Опция (?) не была использована, так как в данном выражении нет ситуаций, где требуется наличие или отсутствие одного символа
\item Операции предпросмотра (lookahead/lookbehind) и другие сложные конструкции не применялись, поскольку они только усложнили бы регулярное выражение в моем случае и не сыграли бы роли.
\item Классы символов не вводились, так как базовые символы достаточно хорошо описывают необходимый язык
\end{itemize}
\end{enumerate}


\end{document}